%\section*{Лекція 12. Основи молекулярно-кінетичної теорії газів}


%% --------------------------------------------------------
\section{120}
%% --------------------------------------------------------

Термодинамічний метод по своїй суті є феноменологічним, тобто ґрунтується на дослідних даних або
теоретичних моделях. Однією з таких моделей і, притому досить успішних, стала молекулярно-кінетична
теорія газів. Розвиток її становлення зобов'язаний працям Даніеля Бернуллі, Жозефа Луї Гей-Люссака,
Джона Дальтона, Джеймса Максвелла, Рудольфа Клаузіуса, Людвіга Больцмана, Ральфа Кроніга -- саме тих
вчених, що внесли неоціненний вклад у розвиток власне термодинаміки. Успіх моделі багато в чому
ґрунтується на твердій експериментальній базі, на відміну від умоглядних моделей будови речовини,
характерних для античності і середньовіччя.

Експериментальні дослідження науки про гази дозволили встановити, що всі гази демонструють
універсальні властивості, що в свою чергу дозволило сподіватися на їх універсальну природу. Нагадаємо
собі ці закони:

\begin{enumerate}
    \item Закон Бойля: $pv = \text{const}$ при постійній температурі;
    \item Закон Шарля: для усіх газів $\alpha = \beta = \frac{1}{273}$ за нормальних умов;
    \item Закон Дальтона: $p = p_1 + p_2 + \dots$;
    \item Закон Джоуля: $U = f(T)$, внутрішня енергія ідеального газу є функцією тільки температури і
    не залежить від об'єму;
    \item Закон Гей-Люссака: об'єми газів, які утворюють хімічні сполуки одна з одною, знаходяться у
    простому співвідношенні один з одним та об'ємом продукту сполуки, якщо останній газоподібний.
    Числа ці цілі і притому невеликі.
\end{enumerate}

До цього слід додати і \emph{закон кратних відношень}, справедливий не тільки для газів, згідно якому
ваги елементів, утворюючих хімічну сполуку з однією і тією ж кількістю деякої речовини, відносяться як
цілі числа. Завдяки цьому закону, Дальтон встановив відносні ваги елементів (відносно водню) і зробив
важливий висновок, вперше кількісно, хоча й опосередковано підтверджуючи гіпотезу про атомістичну
будову речовини:

\begin{quote}
    В хімічних реакціях, атоми реагуючих речовин утворюють більш складні з'єднання -- молекули,
    причому всі молекули хімічно чистої речовини влаштовані однаково.
\end{quote}

\begin{table}[h]
    \centering
    \caption{«Атомні» ваги Дальтона}
    \begin{tblr}{
        colspec = {Q[l, m]Q[c ,m]Q[c,m]},
        column{even}={bg=gray! 10}
        }
        \toprule
        Речовина & Символ & Атомна вага \\
        \midrule
        Водень & H & 1 г \\
        Кисень & O & 5.66 г \\
        Вуглець & C & 4.50 г \\
        Азот & N & 4.00 г \\
        Вода & HO & 6.66 г \\
        Окис азоту & NO & 9.66 г \\
        Окис вуглецю & CO & 10.20 г \\
        \bottomrule
    \end{tblr}
    \label{tab:dalton}
\end{table}

Фактично це означає, що гази, демонструючи у всьому універсальність властивостей, різняться лише
масами молекул, їх утворюючих. Досліди, проведені Гей-Люссаком, доводили, що густина газу пропорційна
його відносній молекулярній масі, і саме $\rho = (4.47 \pm 0.02) \times 10^{-3} \cdot \mu \,
[\text{г/см}^3]$. Фактично, це співвідношення між густиною і відносною молекулярною масою, є кількісне
вираження закону Гей-Люссака. Це могло означати тільки одне: в рівних об'ємах при одних і тих же
тисках і температурі міститься рівна кількість молекул. Це твердження отримало назву гіпотези
Авогадро, висунутої їм в 1811 р., гіпотези у тому сенсі, що напряму перерахувати кількість молекул у
повному обсязі немає можливості. Тим не менш, гіпотеза Авогадро, разом із догадкою Дальтона склали
базис молекулярно-кінетичної теорії газу.

Згідно теорії всі гази складаються з молекул, які знаходяться в безперервному русі. Об'єм молекул
набагато менше, ніж об'єм займаний газом. Дійсно, об'єм рідини, одержуваний при стисненні газу,
набагато менше об'єму газу. Молекули взаємодіють лише при зіткненнях, усю решту часу вони рухаються
прямолінійно і рівномірно. В середньому зіткнення між молекулами і стінками посудини мають бути
пружними. В іншому випадку запас енергії всього газу скоро б вичерпувався, а оскільки для газів вірний
закон Джоуля, то температура газу весь час знижувалась би, чого не відбувається.

%% --------------------------------------------------------
\section{121}
%% --------------------------------------------------------

Подивимось які кількісні результати і наслідки спричиняє дана модель газу.

Спочатку обрахуємо тиск, що чинить газ на стінки посудини. Виділимо елемент площі $d\sigma$ на
поверхні посудини. Кількість молекул, що мають швидкість $\vec{v_i}$ в одиниці об'єму позначимо $n_i$.
Підрахуємо кількість ударів, що спричиняються молекулами цієї групи за час $dt$ на стінку посудини
площею $d\sigma$. Всі ці молекули виявляться в косокутному циліндрі з основою $d\sigma$ і висотою
$v_{ix}dt$, де $v_{ix}$ -- проекція швидкості $\vec{v_i}$ на нормаль до поверхні. Тоді,

\begin{equation}\label{eq:collisions}
    Z_i = d\sigma \cdot v_{ix} \cdot dt \cdot n_i
\end{equation}

Далі, зазвичай вважають, що зіткнення зі стінкою є абсолютно пружні. Досить, проте, щоб вони були в
середньому абсолютно пружні, тобто, щоб кількість відлітаючих в результаті зіткнення зі стінкою
молекул з компонентом швидкості $v_{ix}$ в середньому дорівнювало \eqref{eq:collisions}. Тоді імпульс,
який передається стінці цієї групою молекул буде:

\begin{equation}\label{eq:momentum}
    2mv_{ix} \cdot Z_i
\end{equation}

Повний імпульс, який передається молекулами з усіх швидкісних груп, утворюється підсумовуванням за
всіма швидкостями:

\begin{equation}\label{eq:total_momentum}
    p_x = 2m d\sigma \cdot dt \sum_{v_{ix} > 0} n_i v_{ix}^2
\end{equation}

%% --------------------------------------------------------
\section{123}
%% --------------------------------------------------------

Зіставляючи \eqref{eq:pressure_energy} з рівнянням стану ідеального газу, можемо знайти значення для
$v_{\text{ср.к}}$:

\begin{equation}\label{eq:pressure_ideal_gas}
\frac{1}{3}mn\langle v^{2} \rangle = \frac{\nu RT}{V} = p
\end{equation}

але оскільки $nm = \rho$, то знаходимо, що:

\begin{equation}\label{eq:rms_velocity}
v_{\text{с.кв}} = \sqrt{\frac{3p}{\rho}}
\end{equation}

тобто значення середньоквадратичної швидкості визначається лише макроскопічними параметрами.
Зазначимо, що порядок цієї швидкості такий же, як і порядок швидкості звуку:

\begin{equation}\label{eq:sound_velocity}
v_{\text{ср.кв}} = \sqrt{\frac{3}{\gamma}}c
\end{equation}

Це можна вважати логічним, оскільки передача збурень в звуковій хвилі передається молекулами, які
рухаються з тепловими швидкостями. Для $\text{H}_2$ $v_{\text{ср.кв}} = 183900$ см/сек $= 1.839$
км/сек.


%% --------------------------------------------------------
\section{124}
%% --------------------------------------------------------

Перепишемо \eqref{eq:pressure_energy} у вигляді:

\begin{equation}\label{eq:pressure_kinetic}
P = \frac{2}{3}\frac{1}{2}mn\langle v^{2} \rangle = \frac{2}{3}\frac{\langle \varepsilon_{\text{пост}}
\rangle}{V}
\end{equation}

де $\langle \varepsilon_{\text{пост}} \rangle$ - середнє значення суми кінетичних енергій
поступального руху молекул газу. З \eqref{eq:pressure_kinetic} очевидно, що виконується закон
Дальтона, оскільки $\varepsilon_{\text{пост}} = \sum_{j}E_{j}$, де $j$ нумерує відповідний газ суміші.

Зауважимо, що закон Джоуля виконується автоматично, оскільки енергія газу є сума кінетичної енергії
всіх молекул, а вона не залежить від об'єму\footnote{Наявність обертальних і коливальних ступенів волі
також не дає залежності від об'єму, який займається газом.}. Для того, щоб виконувався закон Бойля
необхідно, щоб величина $\frac{2}{3}\langle \varepsilon_{\text{пост}} \rangle$ була однозначною
функцією температури.

Щоб переконатися в цьому, розглянемо два різні гази розділені поршнем масою $М$. Зіткнення молекули
газу 1 з поршнем описується законами класичної механіки:

\begin{align}
&\text{ЗЗІ: } m_{1}v_{1x} + Mu = m_{1}v'_{1x} + Mu' \label{eq:momentum_conservation} \\
&\text{ЗЗЕ: } \frac{m_{1}}{2}v_{1x}^{2} + \frac{M}{2}u^{2} = \frac{m_{1}}{2}v_{1x}^{'2} +
\frac{M}{2}u'^{2} \label{eq:energy_conservation}
\end{align}

Усереднене значення квадрата проекції швидкості на вісь поршня після зіткнення з поршнем буде:

\begin{equation}\label{eq:velocity_after_collision}
\langle v_{1x}^{'2} \rangle = \frac{4M\langle u^{2} \rangle + (M - m_{1})^{2}\langle v_{1x}^{2}
\rangle}{(M + m_{1})^{2}}
\end{equation}

В стані термодинамічної рівноваги ця величина має дорівнювати $\langle v_{1x}^{2} \rangle$. В іншому
випадку середня кінетична енергія газу 1 змінювалась би з часом. Оскільки $pV \sim \langle
E_{\text{пост}} \rangle$, то поршень прийшов би в рух, або тиск змінювався би з часом. З огляду на
рівність $\langle v_{1x}^{2} \rangle = \langle v_{1x}^{'2} \rangle$ після нескладних перетворень
знайдемо, що:

\begin{equation}\label{eq:energy_equality}
\frac{m_{1}}{2}\langle v_{1x}^{2} \rangle = \frac{M}{2}\langle u^{2} \rangle
\end{equation}

Розмірковуючи таким же чином отримаємо таке ж співвідношення для газу 2. Звідси неважко зрозуміти, що:

\begin{equation}\label{eq:energy_equality_gases}
\frac{1}{2}m_{1}\langle v_{1}^{2} \rangle = \frac{1}{2}m_{2}\langle v_{2}^{2} \rangle
\end{equation}

тобто у стані термодинамічної рівноваги середні кінетичні енергії молекул обох газів однакові. Звідси
витікає закон Авогадро. Зокрема, для однакових об'ємів газів знайдемо, що при однакових $р = р_1 =
р_2$, виконується співвідношення:

\begin{equation}\label{eq:avogadro_law}
p_{1} = \frac{1}{3}\frac{N_{1}}{V}m_{1}\langle v_{1}^{2} \rangle =
\frac{1}{3}\frac{N_{2}}{V}m_{2}\langle v_{2}^{2} \rangle = p_{2}
\end{equation}

Тоді, з огляду на \eqref{eq:energy_equality_gases} отримуємо, що $N_1 = N_2$. Але оскільки закон
Авогадро потребує ще й рівності температур, нам потрібно знову прийняти, щоб середня поступальна
енергія $\langle \varepsilon_{\text{пост}} \rangle = \frac{1}{2}m\langle v^2 \rangle$, що припадає на
одну молекулу газу, має бути монотонною функцією температури.

%% --------------------------------------------------------
\section{125}
%% --------------------------------------------------------

Зручно прийняти за міру температури:

\begin{equation}\label{eq:energy_temperature}
\Theta = \frac{2}{3}\langle \varepsilon_{\text{пост}} \rangle
\end{equation}

Її називають \emph{енергетичною температурою}. Знайдемо її зв'язок із термодинамічною температурою $T$
Кельвіна. Використовуючи \eqref{eq:energy_temperature} знайдемо з \eqref{eq:pressure_kinetic}:

\begin{equation}\label{eq:pV_relation}
pV = N\Theta
\end{equation}

Позначивши $N_A$ -- число молекул газу в одному молі, порівняємо рівняння Клапейрона-Менделєєва з
\eqref{eq:pV_relation} і отримаємо:

\begin{equation}\label{eq:NA_relation}
N_A\Theta = RT
\end{equation}

звідки можна знайти зв'язок між $\Theta$ і $T$. Таким чином:

\begin{equation}\label{eq:boltzmann_constant}
k = \frac{\Theta}{T} = \frac{R}{N_A} = 1.38 \times 10^{-23} \frac{\text{Дж}}{\text{К}}
\end{equation}
де $k$ --- стала Больцмана. Експериментальне значення больцманівської сталої визначається в
експериментах із броунівським рухом. Значення числа Авогадро $N_A = 6.022 \times 10^{23} \text{
моль}^{-1}$ також можна знайти експериментально, зокрема в експерименті Міллікена по визначенню
елементарного заряду. Таким чином температура з точки зору молекулярно-кінетичної теорії це \emph{міра
кінетичної енергії поступального руху}.

%% --------------------------------------------------------
\section{126}
%% --------------------------------------------------------



Скориставшись рівняннями \eqref{eq:125.1} і \eqref{eq:125.3}, знайдемо, що кінетична енергія
поступального руху молекули газу складає:
\begin{equation}
    \left\langle \varepsilon_{\text{пост}} \right\rangle = \frac{3}{2}kT
    \label{eq:126.1}
\end{equation}

Через те, що поступальному руху в тривимірному просторі відповідає три ступені волі, і з огляду на їх
рівноправність, робимо висновок, що на кожну поступальну ступінь волі молекули газу припадає середня
енергія:
\begin{equation}
    \left\langle \varepsilon_{\text{fr}} \right\rangle = \frac{1}{2}kT
    \label{eq:126.2}
\end{equation}

Це співвідношення є частковим випадком \textit{теореми про рівнорозподіл кінетичної енергії за
ступенями волі}, згідно якої у стані теплової рівноваги кожній ступені волі відповідає в середньому
одна і та ж кінетична енергія. Строго ця теорема доводиться в статистичній фізиці. Тут же приведемо
декілька прикладів, які ілюструють теорему.

Вираз \eqref{eq:126.2} отримано для газу, який має лише поступальні ступені волі. Але поршень, з яким
взаємодіяв газ, також складається з молекул. Подивимось скільки в середньому енергії припадає на
молекули поршня. Зауважимо, для початку, що через те, що $m_{1}\langle v_{1}^{2} \rangle =
m_{2}\langle v_{2}^{2} \rangle = M\langle u^{2} \rangle$, можна стверджувати, що одній ступені волі
поршня також відповідає $\frac{1}{2}kT$. Теж саме, очевидно, можна сказати про будь-яке макроскопічне
тіло. Причому, в загальному випадку руху макроскопічного тіла, як цілого, відповідає енергія
$\frac{3}{2}kT$. Розглянемо рух такого тіла. Швидкість центра мас визначається як:
\begin{equation}
    U = \frac{1}{M}\sum_{i} m_{i}u_{i}
    \label{eq:126.3}
\end{equation}

Зводячи до квадрату та помноживши на $\frac{3}{2}$, знайдемо:
\begin{equation}
    \frac{3}{2}MU^{2} = \frac{3}{2}\frac{1}{M}\sum_{i,j} m_{i}m_{j}u_{i}u_{j}
    \label{eq:126.4}
\end{equation}

Усереднюючи та враховуючи, що через хаотичність теплового руху $\langle u_{i}u_{j} \rangle = 0$,
знаходимо:
\begin{equation}
    \frac{3}{2}M\langle U^{2} \rangle = \frac{3}{2}\frac{1}{M}\sum_{i} m_{i}^{2}\langle u_{i}^{2}
    \rangle = \frac{3}{2}kT
    \label{eq:126.5}
\end{equation}

Але
\begin{equation*}
    \sum_{i} m_{i}^{2}\langle u_{i}^{2} \rangle = Nm_{i}^{2}\langle u^{2} \rangle
\end{equation*}
і оскільки $Nm_{i} = M$, а $m_{i} = m$, то:
\begin{equation}
    \frac{3}{2}m\langle u^{2} \rangle = \frac{3}{2}kT
    \label{eq:126.6}
\end{equation}

Тобто і на кожну ступінь волі молекул поршня припадає в середньому $\frac{1}{2}kT$.

Якщо поршень закріплено на пружині, яка задовольняє закону Гука, то поштовхи молекул будуть збуджувати
гармонічні коливання поршня. Оскільки середнє значення кінетичної енергії дорівнює потенційній при
гармонічних коливаннях, то звідси витікає $\langle E_{\text{пот}} \rangle = \frac{1}{2}kT$. Тобто на
кожну коливальну ступінь волі припадає в середньому $2 \cdot \frac{1}{2}kT$. Можна показати, що і
обертанню тіла як цілого довкола деякої осі також відповідає в середньому $\frac{1}{2}kT$ енергії
теплового руху. Дійсно, кінетична енергія деякого тіла обертання:
\begin{equation}
    \frac{1}{2}I\langle \Omega^{2} \rangle = \frac{1}{2I}\sum_{i} m_{i}^{2}\langle r_{i}^{2}u_{i}^{2}
    \rangle = \frac{1}{2}\frac{\langle L^{2} \rangle}{I} = \frac{1}{2}kT\frac{\sum_{i}
    m_{i}r_{i}^{2}}{I} = \frac{1}{2}kT
    \label{eq:126.7}
\end{equation}
де $I = \sum m_{i}r_{i}^{2}$ --- момент інерції відносно осі обертання. Ці приклади показують, що,
дійсно, в рівновазі на кожну ступінь волі одної молекули відводиться $\frac{1}{2}kT$ кінетичної
енергії.


\section*{127. Теорія питомої теплоємності газів}

Власне у виразі \eqref{eq:126.2} вже міститься теорія питомої тепломісткості газу. Дійсно, для
внутрішньої енергії одноатомного газу, який має три поступальні ступені волі, можна записати:
\begin{equation}
    U = \frac{3}{2}N_{A}kT = \frac{3}{2}RT
    \label{eq:127.1}
\end{equation}

Звідки молярна теплоємність:
\begin{align}
    C_{V} &= \frac{3}{2}R \approx 3\ \text{кал/К/моль} \label{eq:127.2} \\
    C_{P} &= C_{V} + R \approx 5\ \text{кал/К/моль} \label{eq:127.3}
\end{align}

Відповідне значення адіабатичної сталої:
\begin{equation}
    \gamma = \frac{5}{3} = 1.667
    \label{eq:127.4}
\end{equation}
що прекрасно узгоджується з експериментом для більшості газів при кімнатних температурах.

Діючи за індукцією, можемо записати, що:
\begin{align}
    C_{V} &= \frac{f}{2}R \label{eq:127.5} \\
    C_{P} &= \frac{f + 2}{2}R \label{eq:127.6} \\
    \gamma &= \frac{f + 2}{f} \label{eq:127.7}
\end{align}
де $f$ --- кількість ступенів волі молекули.

Так, для двохатомної молекули кількість ступенів волі дорівнює п'яти: три поступальні та дві
обертальні, які відповідають обертанню довкола осей $xx'$ та $zz'$. Відповідно, значення:
\begin{align*}
    C_{V} &= \frac{5}{2}R \approx 5\ \text{кал/К/моль} \\
    C_{P} &\approx 7\ \text{кал/К/моль} \\
    \gamma &= \frac{7}{5} = 1.4
\end{align*}
що непогано узгоджується з експериментом (для $H_{2}$ $\gamma = 1.407$, для $N_{2}$ $\gamma = 1.398$,
для $O_{2}$ $\gamma = 1.398$).

Багатоатомні молекули можна розглядати, як тверді тіла, які мають 6 ступенів волі: 3 поступальні та 3
обертальні. Звідси:
\begin{align*}
    C_{V} &= 3R \approx 6\ \text{кал/К/моль} \\
    C_{P} &= 4R\ \text{кал/К/моль} \\
    \gamma &= 1.333
\end{align*}

\begin{figure}[h]
    \centering
%    \includegraphics[width=0.8\textwidth]{Heat_capacity_1.jpg}
    \caption{Графік залежності теплоємності від температури}
    \label{fig:heat_capacity}
\end{figure}

Розвинений підхід виявляється достатньо ефективним при розрахунку тепломісткості твердих і рідких тіл
при не дуже низьких температурах. Так, якщо прийняти за модель кристалу модель досконалої гратки, у
вузлах якої розташовані атоми, то основним механізмом, що акумулює енергію привнесену ззовні, буде
збудження коливань атомів. Тоді, у рівновазі кожній ступені волі відповідає кінетична енергія
$\frac{1}{2}kT$ і така ж потенційна. Таким чином одній коливальній ступені волі відповідає в
середньому:
\begin{equation}
    \langle E_{\text{кол}} \rangle = kT
    \label{eq:127.8}
\end{equation}

Кожний з атомів має три коливальних ступеня волі, тобто на кожний із атомів припадає енергія $3kT$, а
на 1 моль --- $3RT$. Звідси отримуємо закон Дюлонга-Пті, згідно якому добуток питомої тепломісткості
елемента в твердому стані на його атомарну вагу приблизно дорівнює $3R \approx 6$ кал/моль.

Не дивлячись на значні успіхи класичної теорії тепломісткостей, низку принципових моментів вона все ж
таки не пояснює. Зокрема температурну залежність тепломісткості, практично відсутність внеску
електронів металів в її значення, «незатребуваність» деяких ступенів волі (наприклад обертальну
степінь навколо осі $yy'$ або коливальну степінь, як показано на рисунку \ref{fig:heat_capacity}).
Послідовне пояснення цим особливостям дається лише у межах квантомеханічних уявлень. Найбільш суттєвим
моментом є те, що будь-який фінітний рух у квантовій механіці квантується --- енергія набуває лише
дискретні значення. Так, гармонічні коливання мають спектр:
\begin{equation}
    E_{n} = (n + \tfrac{1}{2})\hbar\omega
    \label{eq:127.9}
\end{equation}

За низьких температур при $kT \ll \hbar\omega$ існують лише «нульові» коливання, які не роблять внесок
у тепломісткість. За температур порядку $T_v \approx \hbar\omega/k$ енергія може ефективно
акумулюватися коливальними ступенями волі. Тим самим тепломісткість зростає. Подібний механізм,
відповідальний за зростання тепломісткості, працює і у випадку обертальних ступенів волі. Обертальний
спектр має вигляд:
\begin{equation}
    E_{l} = \frac{\hbar^{2}}{2I} \cdot l(l + 1)
    \label{eq:127.10}
\end{equation}
де $I$ --- момент інерції молекули. При $kT \ll \hbar^{2}/2I$ коливальні ступені волі не збуджені. Їх
внесок позначається за температур вищих за $T_r = \hbar^{2}/2kI$. Так, для водню $T_r \sim 175$ К.

Вказані причини лише в загальних рисах пояснюють спосіб вирішення труднощів класичної теорії.
Зрозуміло проте, що дискретність енергетичних рівнів не сумісна з теоремою про рівномірний розподіл
енергії за ступенями волі принаймні за низьких температур.
